\documentclass[10pt,twoside,a4paper]{report}
%\documentclass[12pt]{book}
\usepackage[dvips]{graphicx}
%\usepackage[leqno,intlimits,sumlimits]{amsmath}
\usepackage[intlimits,sumlimits]{amsmath}
\usepackage[hmargin={1.5cm,1.5cm},vmargin={3cm,3cm}]{geometry}
\usepackage{latexsym,mathrsfs,gensymb,amssymb}
\usepackage{xspace}
\usepackage{upgreek,pifont}
\usepackage{rotating}
\usepackage{supertabular}
\usepackage{overcite}
\usepackage{lscape}
\usepackage{footnpag}
\usepackage{xcolor}
\usepackage{lscape}
\usepackage{pstricks,pst-node,pst-grad}
\usepackage{fancyhdr}
\pagestyle{fancy}
%\usepackage{fontawesome5}
\usepackage{fontawesome} 
\usepackage{wasysym}
\newcommand{\doublespace}{\addtolength{\baselineskip}{0.25\baselineskip}}



\begin{document}
	{\bf{\Large{Lab 1: Identify Entities and Attributes from a Flat File}}}\\
	
    	{\Large{Problem Statement:}} \\
	
	You are given a CSV export from an e-commerce app that mixes order, customer, and product details in each row. Identify the candidate entities and list their key attributes.\\
	
	{\Large{Answer:}}   Input CSV rows or columns
	\texttt{order\_id}, \texttt{order\_date}, \texttt{customer\_name},
	\texttt{customer\_email}, \texttt{product\_id}, \texttt{product\_name},
	\texttt{quantity}, \texttt{price}, \texttt{shipping\_address}\\
	
	\begin{itemize}
		\item Customers
		    \begin{itemize}
		    	\item  \texttt{customer\_id},
		    	\item  \texttt{customer\_name},
		    	 \item \texttt{customer\_email},
		    	 \item \texttt{shipping\_address}
		    \end{itemize}
		    \item Order
		    \begin{itemize}
		    	\item  \texttt{order\_id},
		    	\item  \texttt{order\_date},
		    	\item{\texttt{quantity}},
		    	\item \texttt{shipping\_address}    
	\end{itemize}
	\item Product
	    \begin{itemize}
	    	\item  \texttt{product\_id},
	    	\item  \texttt{producr\_name},
	    	\item \texttt{price},
	    		    \end{itemize}
	
	\end{itemize}
	
	{\bf{\Large{Lab 2: Classify Relationships and Keys}}}\\
	
{\Large{Problem Statement:}} \\
	
	Using entities from Lab 1, state relationship types and propose PK/FK for each table. Include examples where orders contain multiple products.\\
	
	´\Large {Answer:}
\begin{tabular} {|c|c|c|}
	
	\end{tabular}
	
\end{document}
